\documentclass[]{datasheet}

% Load manifest support and apply document metadata
\usepackage{fiddlie-manifest}
\loadmanifest{manifest.yaml}
\applymanifest

\begin{document}
  % Front matter
  \maketitle
  \tableofcontents

  \section{Introduction}

  This document demonstrates using a YAML manifest file to define document
  metadata instead of manually setting \verb|\title{}|, \verb|\author{}|, etc.

  The manifest file contains all document metadata in a structured format,
  making it easier to manage and update.

  \section{Features}

  \subsection{Automatic Metadata}
  The \verb|\applymanifest| command reads the manifest and sets:
  \begin{itemize}
    \item Document title
    \item Short title (for headers)
    \item Author
    \item Date
    \item Revision number
    \item Document ID
    \item Draft watermark (if enabled)
  \end{itemize}

  \subsection{Revision History}
  The \verb|\makerevisionhistory| command generates a revision table
  from the manifest's history section.

  % Auto-generated revision history from manifest
  \makerevisionhistory

  \importantnotice
\end{document}
